% lecture11.tex: Lecture 11: Gas & Ice Giants
% Planetary Systems, Kapteyn Institute

\documentclass[aspectratio=169, 11pt]{beamer}

% ── Theme and macros ────────────────────────────────────────
\usepackage{../common/beamerthemeIPS}
% macros.tex: Shared math macros for IPS lecture slides
% Mirrors definitions in book/_config.yml for consistency
% Uses \providecommand to avoid conflicts with unicode-math

% ── Derivatives ─────────────────────────────────────────────
\providecommand{\dv}[2]{\frac{\mathrm{d} #1}{\mathrm{d} #2}}
\providecommand{\dd}{}\renewcommand{\dd}{\mathrm{d}}
\providecommand{\pdv}[2]{\frac{\partial #1}{\partial #2}}

% ── Solar quantities ────────────────────────────────────────
\newcommand{\Msun}{M_{\odot}}
\newcommand{\Rsun}{R_{\odot}}
\newcommand{\Lsun}{L_{\odot}}

% ── Planetary quantities ────────────────────────────────────
\newcommand{\Mearth}{M_{\oplus}}
\newcommand{\Rearth}{R_{\oplus}}
\newcommand{\Mjup}{M_{\mathrm{J}}}
\newcommand{\Rjup}{R_{\mathrm{J}}}

% ── Constants ───────────────────────────────────────────────
\newcommand{\kB}{k_{\mathrm{B}}}

% ── Media links ─────────────────────────────────────────────
% Clickable button that opens the animated or video version of a
% figure, e.g. \mediabutton{https://...gif}{Play animation}
\providecommand{\mediabutton}[2]{\href{#1}{\beamergotobutton{#2}}}

\graphicspath{{figures/}{../common/}{../../book/11_gas_ice_giants/figures/}}

% ── Metadata ────────────────────────────────────────────────
\title[Lecture 11: Giants]{Lecture 11: Gas \& Ice Giants\\Jupiter, Saturn, Uranus, Neptune}
\subtitle{Planetary Systems}
\author[L11]{Tim Lichtenberg}
\institute{Kapteyn Astronomical Institute\\University of Groningen}
\date{2026}
\titlebgcredit{Jupiter's south pole, JunoCam. Credit: NASA/JPL-Caltech/SwRI/MSSS}

% ════════════════════════════════════════════════════════════
\begin{document}

% ── Title slide ─────────────────────────────────────────────
\begin{frame}[plain,noframenumbering]
  \titlepage
\end{frame}

% ── Learning objectives ─────────────────────────────────────
\begin{frame}{Learning Objectives}
  By the end of this lecture, you will be able to:
  \vskip8pt
  \begin{itemize}
    \item Compare the interiors and atmospheres of the four giant planets
    \item Derive the Roche limit and apply it to Saturn's rings
    \item Explain why the giant-planet moons host oceans, and use the ice giants as analogues of the commonest exoplanets
  \end{itemize}
\end{frame}

% ── Recap ───────────────────────────────────────────────────
\begin{frame}{Recap: Where We Are in the Course}
  \begin{itemize}
    \item \textbf{L9 (Earth \& Venus):} rocky-planet near-twins with divergent surface evolutions
    \item \textbf{L10 (Mercury \& Mars):} rocky-planet limiting cases (dense core vs.\ volatile-rich)
    \item \textbf{Today:} the four giant planets, $>$99.5\% of all planetary mass beyond the Sun
  \end{itemize}
  \vskip8pt
  \textbf{Today:} How do gas and ice giants work, and why do their moons host oceans?
\end{frame}


% ════════════════════════════════════════════════════════════
\sectionimage{family_portrait}{Solar system family portrait, Voyager 1. Credit: NASA/JPL-Caltech}
\section{Gas Giants and Ice Giants: An Overview}
% ════════════════════════════════════════════════════════════

% ── Two classes, four planets ───────────────────────────────
\begin{frame}{Two Classes, Four Planets}
  \begin{itemize}\setlength\itemsep{4pt}
    \item \textbf{Gas giants} (Jupiter, Saturn): H/He envelopes dominate the mass ($318$ and $95\,\Mearth$), together $\sim 92\%$ of all planetary mass
    \item \textbf{Ice giants} (Uranus, Neptune): H/He only $\sim 10$--$20\%$ of the mass ($14.5$ and $17.1\,\Mearth$); the bulk is water, ammonia and methane fluids
    \item No solid surface (the $1$~bar level is the reference); rapid rotation ($\sim 10$~h) makes them oblate by $6$--$10\%$
  \end{itemize}
\end{frame}

% ── Energy budgets ──────────────────────────────────────────
\begin{frame}{Energy Budgets: Internal Heat Excess}
  \begin{itemize}\setlength\itemsep{4pt}
    \item \textbf{Jupiter ($1.7\times$):} slow Kelvin-Helmholtz gravitational contraction
    \item \textbf{Saturn ($1.8\times$):} gravitational energy release from helium rain
    \item \textbf{Neptune ($2.6\times$) vs.\ Uranus ($1.06\times$):} strong internal flux against near solar equilibrium
  \end{itemize}
  \vskip8pt
  \keyresult{Three of the four giants radiate far more than they receive from the Sun; Uranus is the anomalously cold exception.}
\end{frame}

% ── Rapid rotation ──────────────────────────────────────────
\begin{frame}{Why Rapid Rotation Matters}
  \begin{itemize}\setlength\itemsep{4pt}
    \item \textbf{Rapid spin:} Jupiter ($9$~h $56$~min) and Saturn ($10$~h $33$~min) rotate in $\sim 10$~h
    \item \textbf{Zonal banding:} Coriolis forces organise the circulation into alternating jets
    \item \textbf{Gravity harmonics $J_n$:} oblateness and deep flows encode the mass distribution, measured by Juno and Cassini
  \end{itemize}
\end{frame}


% ════════════════════════════════════════════════════════════
\sectionimage{jupiter_south_pole_juno}{Jupiter's south pole, JunoCam. Credit: NASA/JPL-Caltech/SwRI/MSSS}
\section{Jupiter: Interior and Atmosphere}
% ════════════════════════════════════════════════════════════

% ── Jupiter interior and the dilute core ────────────────────
\begin{frame}
  \ipsherokey[0.6]{jupiter_dilute_core_wahl2017}{Density against fractional radius for Jupiter interior models with a compact core (dashed) and a dilute core; metallic hydrogen begins at $\sim 100$~GPa, $r/R_J \approx 0.85$. Wahl et al.\ (2017).}{Juno's gravity harmonics $J_4$ to $J_{10}$ rule out a compact $\sim 10\,\Mearth$ core: $20$--$40\,\Mearth$ of heavy elements spread over the inner $30$--$50\%$ of the radius, mixed during accretion or by a giant impact (Liu et al.\ 2019).}
\end{frame}

% ── Great Red Spot ──────────────────────────────────────────
\begin{frame}
  \ipsherokey[0.66]{jupiter_grs_juno}{Crescent Jupiter and the Great Red Spot, JunoCam (perijove 3, December 2016). Credit: NASA/JPL-Caltech/SwRI/MSSS / R.\ Tkachenko.}{Alternating zonal cloud belts bracket the anticyclonic Great Red Spot, which has shrunk from $\sim 40{,}000$~km in 1900 to $\sim 14{,}000$~km today.}
\end{frame}

% ── Polar cyclones ──────────────────────────────────────────
\begin{frame}{Polar Cyclones}
  \begin{columns}[T]
    \column{0.5\textwidth}
    \includegraphics[width=\linewidth, height=0.44\textheight, keepaspectratio]{jupiter_north_pole_cyclones_juno}\\
    \footnotesize \emph{North pole}: 1 + 8 cyclones, Adriani et al.\ (2018). Credit: NASA/JPL-Caltech/SwRI/ASI/INAF/JIRAM
    \column{0.5\textwidth}
    \includegraphics[width=\linewidth, height=0.44\textheight, keepaspectratio]{jupiter_south_pole_juno}\\
    \footnotesize \emph{South pole}: 1 + 5 cyclones. Credit: NASA/JPL-Caltech/SwRI/MSSS
  \end{columns}
  \vskip4pt
  \keyresult{Stable polygonal arrangements persist across many Juno orbits; polar weather was not seen at high resolution before 2016.}
\end{frame}

% ── Jets, aurorae, and Great Blue Spot ──────────────────────
\begin{frame}{Jets, Aurorae, and the Great Blue Spot}
  \begin{itemize}\setlength\itemsep{4pt}
    \item \textbf{Jet penetration:} winds extend to $\sim 3000$~km before magnetic braking enforces rigid rotation (Kaspi et al.\ 2018)
    \item \textbf{Io mass loading:} $\sim 1\ \mathrm{t\,s^{-1}}$ of sulfur and oxygen ions feed the strongest aurorae in the solar system; UV footprints mark the field lines to Io, Europa and Ganymede
    \item \textbf{Great Blue Spot:} isolated equatorial magnetic anomaly mapped by Juno (Connerney et al.\ 2022)
  \end{itemize}
\end{frame}


% ════════════════════════════════════════════════════════════
\sectionimage{galilean_moons}{The Galilean moons at correct relative size: Io, Europa, Ganymede, Callisto. Credit: NASA/JPL/DLR}
\section{The Galilean Moons}
% ════════════════════════════════════════════════════════════

% ── Io ──────────────────────────────────────────────────────
\begin{frame}
  \ipsherokey[0.65]{io_loki_volcano}{Loki Patera lava lake on Io, Voyager 1 (1979). Credit: NASA/JPL.}{About $400$ active volcanic centres and $\sim 10^{14}$~W of tidal heat driven by the 1:2:4 Laplace resonance, predicted by Peale et al.\ (1979) days before Voyager 1; the surface is resurfaced on $\sim$Myr timescales.}
\end{frame}

\begin{frame}
  \ipshero{io_tvashtar_eruption}{The Tvashtar Catena fire fountain on Io in a Galileo composite of orbits C21, I25 (26 November 1999) and I27 (22 February 2000). Credit: NASA/JPL/University of Arizona.}
\end{frame}

\begin{frame}
  \ipsherokey[0.56]{io_tidal_park2024}{Io's measured tidal response, $|k_2|/Q$ against $\mathrm{Re}(k_2)$, for interior models without (a) and with (b) a global magma ocean; green boxes are the Juno constraints. Reproduced from Park et al.\ (2025).}{Models without a magma ocean pass through the Juno box, while a magma ocean shallower than about 500~km gives a far larger $\mathrm{Re}(k_2)$ than measured: the data favour a mostly solid mantle.}
\end{frame}

% ── Europa ──────────────────────────────────────────────────
\begin{frame}
  \ipsherokey{europa_galileo_mosaic}{Europa surface mosaic from the Galileo orbiter. Credit: NASA/JPL-Caltech/SETI Institute.}{Global lineae fractures and disrupted chaos terrain reflect ongoing tidal heating; the surface is only $40$--$90$~Myr old.}
\end{frame}

\begin{frame}
  \ipsherokey[0.66]{europa_chaos_terrain}{Conamara Chaos on Europa, Galileo (PIA01640). Credit: NASA/JPL-Caltech/ASU.}{Chaos morphology records episodic basal melting, and the induced dipole seen by the Galileo magnetometer requires a global conducting ocean (Khurana et al.\ 1998): $6$--$25$~km of ice over a $\sim 100$~km ocean, about one Earth-ocean volume.}
\end{frame}

% ── Ganymede ────────────────────────────────────────────────
\begin{frame}
  \ipsherokey[0.62]{ganymede_juno_closeup}{Ganymede imaged by JunoCam during the 7 June 2021 close flyby. Credit: NASA/JPL-Caltech/SwRI/MSSS.}{Radius $2634$~km, larger than Mercury; the only moon with an intrinsic dynamo (Galileo 1996); fully differentiated, with a $\sim 100$~km ocean below a $\sim 150$~km ice shell (Saur et al.\ 2015).}
\end{frame}

\begin{frame}
  \ipshero[0.76]{ganymede_grooves}{Grooved bright terrain (Lagash Sulcus) cutting ancient cratered dark terrain in Marius Regio on Ganymede, Galileo, 6 June 1997, about 288 m per pixel. Credit: NASA/JPL-Caltech/Brown University.}
\end{frame}

% ── Callisto ────────────────────────────────────────────────
\begin{frame}
  \ipsherokey[0.61]{callisto_global}{Callisto in global colour view, assembled from Galileo and Voyager images. Credit: NASA/JPL-Caltech.}{$R = 2410$~km, density $1834\ \mathrm{kg\,m^{-3}}$, $C/MR^2 \approx 0.355$: only partially differentiated (Anderson et al.\ 2001), outside the Laplace resonance and without tectonic resurfacing, yet an induced field points to a subsurface ocean.}
\end{frame}

\begin{frame}
  \ipshero{callisto_cutaway}{Interior structures of the four Galilean moons at comparable scale: only Callisto (right) keeps an undifferentiated rock-ice interior. Credit: NASA/JPL-Caltech.}
\end{frame}


% ════════════════════════════════════════════════════════════
\sectionimage{saturn_rings}{Saturn and its rings, Cassini (PIA21046). Credit: NASA/JPL-Caltech/SSI}
\section{Saturn: Interior and Rotation}
% ════════════════════════════════════════════════════════════

% ── Saturn interior: helium rain ────────────────────────────
\begin{frame}{Saturn Interior: Lower Pressures, Helium Rain}
  \begin{itemize}\setlength\itemsep{4pt}
    \item Mass $95\,\Mearth$, central pressure $\sim 1/3$ of Jupiter's: the molecular-to-metallic hydrogen transition sits deeper, at milder conditions
    \item Helium becomes immiscible in metallic H at $1$--$3$~Mbar, $5$--$10$~kK; the falling droplets release gravitational energy and close Saturn's heat budget
    \item The rain depletes the upper envelope in helium, consistent with the Voyager observations; the Cassini Grand Finale (2017) mapped the gravity field
  \end{itemize}
\end{frame}

\begin{frame}
  \ipsherokey[0.65]{saturn_kronoseismology_mankovich2021}{Kronoseismology: f-mode resonances of Saturn's interior imprinted as waves in the C ring. Mankovich \& Fuller (2021).}{The ring waves require a stably stratified dilute core out to $\sim 60\%$ of $R_S$ holding $\sim 17\,\Mearth$ of heavy elements.}
\end{frame}

\begin{frame}
  \ipshero[0.76]{saturn_interior_mankovich2021}{Saturn's heavy-element fraction $Z(r)$ (top), density (middle) and Brunt-Vaisala frequency (bottom) against fractional radius, coloured by model likelihood, maximum-likelihood profile in yellow. Reproduced from Mankovich \& Fuller (2021).}
\end{frame}

% ── Saturn rotation from ring seismology ────────────────────
\begin{frame}
  \ipsherokey[0.61]{saturn_rotation_mankovich2019}{RMS pattern-speed residual between interior models and the observed C-ring density waves against the assumed rotation period. Mankovich et al.\ (2019).}{Saturn's nearly axisymmetric dipole gives no radio period; the C-ring f-mode frequencies give $P_S = 10$~h $33$~min $38$~s and resolve the decades-long Voyager-Cassini ambiguity.}
\end{frame}

% ── Saturn atmosphere and hexagon ───────────────────────────
\begin{frame}
  \ipsherokey[0.61]{saturn_hexagon_jet}{Saturn's hexagonal polar jet, Cassini false-colour view (PIA14946). Credit: NASA/JPL-Caltech/SSI/Hampton University.}{The same $\mathrm{NH_3}$/$\mathrm{NH_4SH}$/$\mathrm{H_2O}$ cloud decks as Jupiter, deeper and muted; the hexagon at $\sim 78^\circ$\,N is a six-sided standing Rossby wave seen since Voyager (1980 to 1981), and the equatorial jet reaches $\sim 400\ \mathrm{m\,s^{-1}}$.}
\end{frame}


% ════════════════════════════════════════════════════════════
\sectionimage{saturn_rings}{Saturn's main ring system, Cassini (PIA21046). Credit: NASA/JPL-Caltech/SSI}
\section{Saturn's Rings}
% ════════════════════════════════════════════════════════════

\begin{frame}
  \ipsherokey{saturn_cassini_division}{Natural-colour radial scan across the C ring, B ring, Cassini Division and A ring, Cassini (PIA08389). Credit: NASA/JPL-Caltech/SSI.}{The rings span $67{,}000$--$140{,}000$~km, consist of $>95\%$ pure water ice, and are only $\sim 10$~m thick.}
\end{frame}

% ── Embedded moons and propellers ───────────────────────────
\begin{frame}
  \ipsherokey[0.61]{saturn_propeller_targeted}{Two Cassini close-ups of the same propeller feature in the A ring, the wake of an embedded $\sim 1$~km moonlet. Credit: NASA/JPL-Caltech/Space Science Institute.}{Prometheus and Pandora shepherd the F ring, Pan clears the Encke Gap, and the propeller moonlets show orbital migration directly within a disk (Tiscareno et al.\ 2013).}
\end{frame}

% ── Ring youth and origin ───────────────────────────────────
\begin{frame}
  \ipsherokey[0.65]{saturn_rings}{Saturn and its rings, Cassini. Credit: NASA/JPL-Caltech/SSI.}{Only $\sim 1.5 \times 10^{19}$~kg, $\sim 40\%$ of Mimas (Iess et al.\ 2019), and still bright despite micrometeoroid infall: about $100$~Myr old (Crida et al.\ 2019), draining as ring rain on $10^8$~yr scales, formed by the tidal disruption of an icy moon (Wisdom et al.\ 2022).}
\end{frame}


% ════════════════════════════════════════════════════════════
\sectionimage{family_portrait}{Solar system family portrait. Credit: NASA/JPL-Caltech}
\section{Blackboard Derivation: The Roche Limit}
% ════════════════════════════════════════════════════════════

% ── Goal and strategy ───────────────────────────────────────
\begin{frame}{Goal and Strategy}
  \begin{columns}[c]
    \column{\recapfigcolnarrow}
    \ipsrecapfig[0.62\textheight]{board_sketch_l11}
    \column{\recaptextcolwide}
    \small
    \textbf{Goal:} the distance inside which tidal stretch beats self-gravity (Roche 1849): rings persist, moons cannot accrete.
    \vskip4pt
    \begin{itemize}\setlength\itemsep{2pt}
      \item Planet: mass $M_p$, radius $R_p$, mean density $\rho_p$
      \item Satellite: mass $M_s$, radius $R_s$, density $\rho_s$, at distance $d \gg R_s$
      \item Balance the tidal stretch $\Delta a_{\mathrm{tidal}} \approx 2 G M_p R_s / d^3$ against the surface self-gravity $G M_s / R_s^2$
    \end{itemize}
    \vskip4pt
    \keyresult{To the board.}
  \end{columns}
\end{frame}

% ── Fluid vs rigid Roche limit ──────────────────────────────
\begin{frame}{Fluid vs.\ Rigid Roche Limit}
  \begin{itemize}
    \item Rigid case: the body holds its spherical shape until disruption: $d_R^{\mathrm{rigid}} = 2^{1/3} R_p (\rho_p/\rho_s)^{1/3} \approx 1.26\, R_p (\rho_p/\rho_s)^{1/3}$
    \item Fluid case: the body deforms into an ellipsoid first and disrupts further out; Roche's calculation gives the prefactor $\approx 2.46$:
  \end{itemize}
  \[
    \boxed{\; d_R \approx 2.46\, R_p \left(\frac{\rho_p}{\rho_s}\right)^{1/3} \;}
  \]
  Real bodies with material strength sit between the rigid and fluid limits.
\end{frame}

\begin{frame}
  \ipshero{roche_geometry}{Roche limit geometry: (a) a fluid satellite elongates before it disrupts at $d_R$; (b) differential tidal acceleration against surface self-gravity for an icy body near Saturn. Course figure.}
\end{frame}

% ── Application: Saturn's rings ─────────────────────────────
\begin{frame}{Apply to Saturn's Rings}
  \[
    R_p = 58{,}232~\mathrm{km}, \quad \rho_p = 687~\mathrm{kg\,m^{-3}}, \quad \rho_s \approx 1000~\mathrm{kg\,m^{-3}}
  \]
  \[
    d_R \approx 2.46 \times 58{,}232 \times (0.687)^{1/3}~\mathrm{km}
        \approx 2.46 \times 58{,}232 \times 0.883
  \]
  \[
    \boxed{\; d_R \approx 126{,}000~\mathrm{km} \;}
  \]
  \vskip6pt
  \keyresult{Within $\sim 10\%$ of A-ring outer edge ($\sim 137{,}000$~km). Discrepancy reflects particle porosity, ice rigidity, and resonant structures.}
\end{frame}


% ── Break ───────────────────────────────────────────────────
\breakslide


% ════════════════════════════════════════════════════════════
\sectionimage{titan_lakes_cassini}{Titan's polar lakes, Cassini RADAR. NASA/JPL-Caltech/ASI/Cornell}
\section{Saturn's Moons: Titan, Enceladus, and More}
% ════════════════════════════════════════════════════════════

% ── Titan ───────────────────────────────────────────────────
\begin{frame}
  \ipsherokey[0.61]{huygens_titan_descent}{Huygens descent imagery of Titan's surface, 14 January 2005 (PIA07870). Credit: ESA/NASA/JPL-Caltech/University of Arizona.}{$R = 2575$~km, the second largest moon; an $\mathrm{N_2}$ atmosphere at $1.5$~bar ($4\times$ Earth's column density) and a $94$~K surface near methane's triple point drive a methane cycle of clouds, rain, rivers and polar hydrocarbon lakes.}
\end{frame}

\begin{frame}
  \ipshero{titan_lakes_cassini}{Cassini radar map of Titan's polar hydrocarbon seas, including Ligeia Mare and Kraken Mare, dark and radar-smooth. Credit: NASA/JPL-Caltech/ASI/Cornell.}
\end{frame}

% ── Enceladus ───────────────────────────────────────────────
\begin{frame}
  \ipsherokey[0.66]{enceladus_geyser_basin}{Cryovolcanic geyser jets erupting from the south-polar tiger stripe fractures of Enceladus, Cassini. Credit: NASA/JPL-Caltech/SSI.}{Dozens of jets powered by $>10$~GW of tidal heat spray subsurface ocean water directly into space: the most accessible ocean world.}
\end{frame}

\begin{frame}{Enceladus: Habitability Ingredients}
  \begin{itemize}\setlength\itemsep{4pt}
    \item Global subsurface ocean beneath a $20$--$30$~km ice shell; the 2:1 resonance with Dione forces the eccentricity that powers the tidal heating
    \item Plume $\mathrm{H_2}$ (Waite et al.\ 2017) and silica nanoparticles (Hsu et al.\ 2015): chemical disequilibrium from active rock-water reactions at hydrothermal vents
    \item Phosphates (Postberg et al.\ 2023): an essential biochemical building block
  \end{itemize}
\end{frame}

\begin{frame}
  \ipsherokey[0.66]{enceladus_tiger_stripes}{The four parallel tiger stripe fractures near the south pole of Enceladus that feed the plumes. Credit: NASA/JPL-Caltech/SSI.}{The fractures are warmer than the surrounding terrain by tens of kelvins, and their ages and orientations track the stress field of the eccentric orbit around Saturn.}
\end{frame}

\begin{frame}
  \ipshero{enceladus_tiger_thermal}{Cassini CIRS thermal map of the south polar terrain of Enceladus: the endogenic heat flow concentrates along the four tiger stripes. Credit: NASA/JPL-Caltech/GSFC/SwRI/SSI.}
\end{frame}

% ── Mid-sized moons ─────────────────────────────────────────
\begin{frame}{Mid-Sized Moons: Mimas and Iapetus}
  \begin{columns}[T]
    \column{0.5\textwidth}
    \includegraphics[width=\linewidth, height=0.55\textheight, keepaspectratio]{mimas_close}\\
    \footnotesize \emph{Mimas, Herschel crater} (near the disruption threshold). Credit: NASA/JPL-Caltech/Space Science Institute
    \column{0.5\textwidth}
    \includegraphics[width=\linewidth, height=0.55\textheight, keepaspectratio]{iapetus_bright_dark}\\
    \footnotesize \emph{Iapetus, two-toned hemispheres}. Credit: NASA/JPL-Caltech/Space Science Institute
  \end{columns}
\end{frame}


% ════════════════════════════════════════════════════════════
\sectionimage{uranus_voyager}{Uranus, Voyager 2 (1986). NASA/JPL-Caltech}
\section{Ice Giants: Composition and Interior}
% ════════════════════════════════════════════════════════════

% ── Ice giants overview ─────────────────────────────────────
\begin{frame}{Ice Giants: The Exotic Twins}
  \begin{itemize}\setlength\itemsep{4pt}
    \item Uranus: $14.5\,\Mearth$, $4.0\,\Rearth$; Neptune: $17.1\,\Mearth$, $3.9\,\Rearth$
    \item H/He envelope only $10$--$20\%$ of the mass; the bulk is water, ammonia and methane in a compressed fluid state
    \item One spacecraft visit each: Voyager 2 in January 1986 and August 1989
  \end{itemize}
  \vskip8pt
  \keyresult{The most under-explored major planets: the data are $35$--$40$~yr old.}
\end{frame}

% ── Possible interior structures ────────────────────────────
\begin{frame}
  \ipsherokey{ice_giant_structures_helled2020}{Possible ice-giant interiors, from sharp layer boundaries (a) to a fully gradual composition gradient (d). Helled et al.\ (2020).}{The Voyager $J_2$ and $J_4$ cannot break the degeneracy between these density profiles.}
\end{frame}

\begin{frame}
  \ipshero[0.77]{ice_giant_density_helled2020}{Density against radius for Uranus (blue) and Neptune (black): empirical profiles (solid) and three-layer models (dashed) match the gravity data equally well. Reproduced from Helled et al.\ (2020).}
\end{frame}

% ── Superionic ice and magnetic fields ──────────────────────
\begin{frame}{Superionic Ice and Tilted Dynamos}
  \begin{itemize}\setlength\itemsep{4pt}
    \item \textbf{Superionic ice:} at $100$--$400$~GPa a solid oxygen lattice with mobile protons conducts (Millot et al.\ 2019)
    \item \textbf{Tilted, offset fields:} $59^\circ$ (Uranus) and $47^\circ$ (Neptune), offset $\sim 1/3\,R_p$ and $\sim 1/2\,R_p$ from the centre
    \item \textbf{Thin-shell dynamos:} convection in an outer shell explains the multipolar geometry
  \end{itemize}
  \vskip6pt
  \keyresult{Superionic ice powers dynamos without metallic hydrogen.}
\end{frame}

\begin{frame}
  \ipshero[0.76]{soderlund2020_planetary_field_morphologies}{Surface radial magnetic field of Mercury, Earth, Jupiter, Saturn, Uranus and Neptune, each on its own colour scale: aligned dipoles at Jupiter and Saturn, tilted and offset multipoles at Uranus and Neptune. Reproduced from Soderlund \& Stanley (2020).}
\end{frame}


% ════════════════════════════════════════════════════════════
\sectionimage{uranus_clouds_voyager}{Uranus clouds, Voyager 2 (PIA18182)}
\section{Uranus: The Tilted Planet}
% ════════════════════════════════════════════════════════════

\begin{frame}
  \ipsherokey{uranus_impact_kegerreis2018}{SPH simulation of a giant impact on proto-Uranus. Kegerreis et al.\ (2018).}{An oblique collision tilted the spin axis to $97.8^\circ$, giving $42$-yr polar seasons, and shed an equatorial debris disk.}
\end{frame}

% ── Uranus: Voyager 2 observations ──────────────────────────
\begin{frame}{Uranus: Featureless at Voyager, Active Now}
  \begin{columns}[T]
    \column{0.5\textwidth}
    \includegraphics[width=\linewidth, height=0.55\textheight, keepaspectratio]{uranus_voyager}\\
    \footnotesize Uranus, Voyager 2 (1986). Credit: NASA/JPL-Caltech
    \column{0.5\textwidth}
    \includegraphics[width=\linewidth, height=0.55\textheight, keepaspectratio]{uranus_clouds_voyager}\\
    \footnotesize Enhanced cloud bands, Voyager 2 (PIA18182). Credit: NASA/JPL-Caltech
  \end{columns}
  \vskip4pt
  Voyager 2 (1986) observed Uranus near solstice: pole-on orientation and muted cloud bands.
\end{frame}

% ── Modern Uranus ───────────────────────────────────────────
\begin{frame}
  \ipsherokey{uranus_cyclone}{Uranus with its north-polar cyclone in recent observations.}{Since the 2007 equinox, toward northern summer (2030), cloud activity and storms increase: a 2014 storm system from the ground, vivid rings and a polar cap from JWST (2023), and a north-polar cyclone confirmed by the VLA and JWST.}
\end{frame}

% ── Uranus heat budget anomaly ──────────────────────────────
\begin{frame}{Uranus: Why So Cold Inside?}
  \begin{itemize}\setlength\itemsep{4pt}
    \item \textbf{Heat anomaly:} $\sim 1.06\times$ the absorbed solar flux (Pearl et al.\ 1990); Neptune $\sim 2.6\times$
    \item \textbf{Trapped heat:} a giant impact may have left a stably stratified interior that suppresses convection
    \item \textbf{Latent heat:} deep condensation of volatiles could delay the cooling
  \end{itemize}
  \vskip8pt
  \keyresult{Uranus emits almost no heat beyond the sunlight it absorbs; Neptune does.}
\end{frame}


% ════════════════════════════════════════════════════════════
\sectionimage{neptune_great_dark_spot}{Neptune Great Dark Spot, Voyager 2 (1989). NASA/JPL-Caltech}
\section{Neptune and Triton}
% ════════════════════════════════════════════════════════════

\begin{frame}
  \ipsherokey{neptune_great_dark_spot}{Neptune's Great Dark Spot, Voyager 2 (1989). Credit: NASA/JPL-Caltech.}{An anticyclonic storm the size of Earth, bordered by bright methane cirrus, seen during the 1989 flyby and gone by 1994 (HST); new dark spots emerge periodically.}
\end{frame}

\begin{frame}
  \ipsherokey[0.62]{neptune_scooter}{The bright ``Scooter'' cloud feature near $42^\circ$S on Neptune, Voyager 2 (1989). Credit: NASA/JPL-Caltech.}{The Scooter drifts westward relative to the $16.1$~h interior rotation and the equatorial jet peaks at $\sim 400\ \mathrm{m\,s^{-1}}$ on $1/900$ of Earth's insolation: the winds draw on the $\sim 2.6\times$ internal heat excess (Pearl \& Conrath 1991).}
\end{frame}

\begin{frame}
  \ipsherokey[0.61]{triton_map}{Triton's southern hemisphere, Voyager 2 mosaic, with cantaloupe terrain and dark geyser streaks. Credit: NASA/JPL-Caltech.}{A captured Kuiper Belt object: regular moons form prograde, so the retrograde orbit ($i \approx 157^\circ$) means capture; active $8$-km $\mathrm{N_2}$ geysers and a surface younger than $\lesssim 100$~Myr.}
\end{frame}


% ════════════════════════════════════════════════════════════
\sectionimage{neptune_rings_voyager}{Neptune ring system, Voyager 2 (1989). NASA/JPL-Caltech}
\section{Ice Giant Rings}
% ════════════════════════════════════════════════════════════

% ── Ice giant ring systems ──────────────────────────────────
\begin{frame}{Ice Giant Rings: Narrow, Dark, and Carbonaceous}
  \begin{itemize}\setlength\itemsep{4pt}
    \item \textbf{Occultation discovery:} Uranus's rings were found in 1977 by stellar occultation (13 narrow rings known today)
    \item \textbf{Compositional contrast:} Saturn's rings are $>95\%$ bright water ice, ice-giant rings are dark organics
    \item \textbf{Disruption origin:} collisional destruction of small inner moonlets rather than a single large icy parent
  \end{itemize}
\end{frame}

% ── Neptune rings: arcs ─────────────────────────────────────
\begin{frame}{Neptune Rings: Arcs Held by Resonance}
  \begin{columns}[T]
    \column{0.5\textwidth}
    \includegraphics[width=\linewidth, height=0.5\textheight, keepaspectratio]{neptune_rings_voyager}\\
    \footnotesize Adams ring + Le Verrier ring. Credit: NASA/JPL-Caltech
    \column{0.5\textwidth}
    \includegraphics[width=\linewidth, height=0.5\textheight, keepaspectratio]{neptune_rings_voyager_arc}\\
    \footnotesize Long-exposure forward-scattered view. Credit: NASA/JPL-Caltech
  \end{columns}
  \vskip4pt
  \keyresult{The Adams ring holds 5 distinct \emph{arcs}, trapped by a mean-motion resonance with Galatea.}
\end{frame}


% ════════════════════════════════════════════════════════════
\sectionimage{europa_clipper_concept}{Europa Clipper concept. NASA/JPL-Caltech}
\section{Comparative Payoff and Exploration Frontier}
% ════════════════════════════════════════════════════════════

\begin{frame}{Why Gas Giants and Ice Giants Diverged}
  \begin{itemize}\setlength\itemsep{4pt}
    \item \textbf{Timing:} Jupiter and Saturn reached runaway gas accretion while the nebula was still massive; Uranus and Neptune grew too slowly and caught only $1$--$4\,\Mearth$ of H/He
    \item \textbf{Result:} two H/He-dominated planets with metallic hydrogen dynamos, two ice-dominated planets with superionic-ice dynamos
    \item \textbf{Common threads:} rapid rotation, internal heat (except Uranus), rings and regular moon systems, one family of physics
  \end{itemize}
\end{frame}

\begin{frame}
  \ipsherokey[0.65]{helled2014_core_accretion}{Giant-planet mass evolution by core accretion at 5.2~AU: core mass (solid), envelope mass (dash-dotted) and total mass (dashed). Reproduced from Helled et al.\ (2014).}{Rapid core assembly, a slow quasi-hydrostatic envelope contraction, then runaway gas accretion once the envelope mass exceeds the core mass.}
\end{frame}

\begin{frame}
  \ipsherokey[0.67]{europa_clipper_concept}{Europa Clipper at Europa, artist's concept. Credit: NASA/JPL-Caltech.}{Europa Clipper (launched 2024, $\sim 50$ Europa flybys) and JUICE (launched 2023, Ganymede orbit 2034) are en route to test the habitability of the Jovian ocean worlds.}
\end{frame}


% ════════════════════════════════════════════════════════════
\section*{Summary}
% ════════════════════════════════════════════════════════════

% ── Summary ────────────────────────────────────────────────
\begin{frame}{Summary: Today's Course-Level Takeaways}
  \begin{enumerate}\setlength\itemsep{6pt}
    \item \textbf{Two classes, dilute cores:} gas giants (H/He-dominated, Saturn $\sim 80\%$ H/He by mass) against ice giants (H/He only $10$--$20\%$); Juno and ring seismology reveal diffuse cores in Jupiter and Saturn
    \item \textbf{Roche limit:} $d_R \approx 2.46\,R_p(\rho_p/\rho_s)^{1/3}$ bounds rings and moon accretion; tidal heating maintains the oceans of Europa and Enceladus
    \item \textbf{Under-explored ice giants:} tilted dynamos powered by superionic ice, and the closest analogues of the sub-Neptunes of Lecture 13
  \end{enumerate}
  \vskip8pt
  \textbf{Next: Lecture 12}
\end{frame}

% ── Final slide ─────────────────────────────────────────────
\begin{frame}[plain,noframenumbering]
  \begin{tikzpicture}[remember picture, overlay]
    \ipscontain{title_background}
    \fill[ipsPrimary, opacity=0.70] (current page.north west) rectangle (current page.south east);
  \end{tikzpicture}
  \vfill
  \begin{center}
    {\color{white}\Huge \textbf{Questions?}}\\[14pt]
    {\color{white}\large Next: \textbf{Lecture 12. Small Bodies: Meteorites, Asteroids, Comets}}\\[10pt]
    {\color{white}\normalsize Lecture notes: \texttt{formingworlds.github.io/IntroductionToPlanetaryScience}}
  \end{center}
  \vfill
\end{frame}

\end{document}
